\documentclass[journal=jacsat,manuscript=article]{achemso}
\usepackage[T1]{fontenc}
\author{Haruka Shimotsuki}
\affiliation[Kitamachi University]{Department of Chemistry, Kitamachi University, Kitamachi 100-0000, Japan}
\author{Marcus B. Ellery}
\affiliation[Westbrook Institute]{Department of Chemistry, Westbrook Institute of Technology, Westbrook 00000, Freedonia}
\author{Yuki Tanabe}
\email{tanabe@example.org}
\affiliation[Kitamachi University]{Department of Chemistry, Kitamachi University, Kitamachi 100-0000, Japan}
\title{Room-Temperature Assembly of a Fictional Porous Framework for Selective Water Harvesting}
\keywords{metal--organic frameworks, water harvesting, porous materials}
\begin{document}
\begin{abstract}
Atmospheric water harvesting requires porous materials that capture water at low humidity and release it with little energy.
Here we report a fictional zirconium framework, KU-42, that assembles at room temperature in water within ten minutes.
KU-42 shows a steep water uptake step at 18\% relative humidity and retains its capacity over 500 adsorption--desorption cycles.
In situ diffraction reveals that water molecules first bind to open metal sites and then fill the pores cooperatively, which explains the sharp uptake.
A prototype device delivers 1.2~L of water per kilogram of material per day under desert-like conditions.
These results show that mild synthesis and high performance are not mutually exclusive for water-harvesting frameworks.
\end{abstract}
\section{Introduction}
Water scarcity affects billions of people, and harvesting water from air is a promising complement to desalination~\cite{shimotsuki2021,ellery2022}.
Porous frameworks are attractive sorbents because their pore size and chemistry can be tuned at the molecular level.
However, most high-performing frameworks require solvothermal synthesis at elevated temperatures, which limits their scalability~\cite{tanabe2023}.

In this work we show that a carefully chosen linker allows a zirconium framework to assemble in water at room temperature.
We characterize its structure, its water sorption behaviour and its stability, and we demonstrate its use in a simple harvesting device.
The approach relies only on commercially available reagents and could be adapted to related frameworks.

\section{Results and Discussion}
The framework was obtained as a white powder.
\bibliography{refs}
\end{document}
